% TOP_PROBLEM: {"id":30007583,"problem_number":"LOCAL-30007583","title":"A portable model of cognitive biases","category_id":21,"category":{"id":21,"name":"economics","display_name":"Economics","description":"Open economic theory statements, published research agendas, and proposed scoped specifications; question provenance and historical priority are qualified per record.","slug":"economics","order_index":21,"created_at":"2026-10-08T00:00:00Z"},"status":"research_question","external_url":null,"legacy_ids":[],"published":false,"statement":"Identify whether the specified cognitive-uncertainty model predicts choice biases with portable parameters across economic decision tasks.","economics":{"original_id":72,"field":"Behavioral and experimental economics","kind":"empirical","formalization_basis":"adapted_research_question","source_status":"research_agenda","priority_status":"earliest_found","priority_note":"The 2023 discussion is the earliest full-text statement verified here. A 2019 working paper exists, but its priority for these particular open extensions was not verified.","earliest_verified_reference":{"authors":["Benjamin Enke","Thomas Graeber"],"title":"Cognitive Uncertainty","year":2023,"url":"https://thomasgraeber.com/assets/papers/EnkeGraeber_CognitiveUncertainty.pdf","doi":"10.1093/qje/qjad025","locator":"Author manuscript dated March 14, 2023, §8 Discussion, printed pp. 33–35 (PDF pp. 34–36).","evidence_summary":"The discussion identifies missing general theories of cognitive defaults and task complexity, asks when cognitive uncertainty reflects actual noise, and suggests applications in additional decision domains."},"current_reference":{"authors":["Benjamin Enke","Thomas Graeber"],"title":"Cognitive Uncertainty","year":2023,"url":"https://thomasgraeber.com/assets/papers/EnkeGraeber_CognitiveUncertainty.pdf","doi":"10.1093/qje/qjad025","locator":"Author manuscript dated March 14, 2023, §8 Discussion, printed pp. 33–35 (PDF pp. 34–36).","evidence_summary":"The discussion identifies missing general theories of cognitive defaults and task complexity, asks when cognitive uncertainty reflects actual noise, and suggests applications in additional decision domains."},"formalization_note":"The source explicitly discusses missing theories, calibration, and additional domains. The unified prediction equation, held-out evaluation, and universal-model wording are our adaptation."},"statusQualification":"Published question or research agenda verified at the cited locator. The mathematical specification is adapted for this catalogue and includes additional assumptions; source publication does not establish first-ever priority or certify this exact model remains unsolved. The source explicitly discusses missing theories, calibration, and additional domains. The unified prediction equation, held-out evaluation, and universal-model wording are our adaptation.","statusReviewedAt":"2026-10-08"}
% FIRST_POSTED: 2026-10-08
% ENTRY_KIND: statement_only
% !TeX program = lualatex
\documentclass[11pt]{article}
\usepackage{fontspec}
\usepackage[a4paper,margin=25mm]{geometry}
\usepackage{amsmath,amssymb,mathtools}
\usepackage{xurl}
\usepackage[hidelinks]{hyperref}
\newcommand{\sref}[2]{\hyperlink{source:#1:#2}{[S#2]}}
\setlength{\emergencystretch}{3em}
\begin{document}
% problemId:30007583
\section*{A portable model of cognitive biases}\label{30007583}
\subsection{Definitions and mathematical statement}
Let \(d\) index risky choice, ambiguous choice, belief updating, and intertemporal choice. In each domain individual \(i\) faces a task with observed attributes \(X_{id}\), gives response \(Y_{id}\), and reports cognitive uncertainty \(Q_{id}\in[0,1]\), defined as subjective doubt that the chosen action maximizes expected utility. Let \(m_d(X_{id};\theta_i)\) denote a domain-specific benchmark decision, where \(\theta_i\) contains stable preference or belief parameters. A candidate shared mechanism has the form
\[Y_{id}=(1-\lambda_{id})m_d(X_{id};\theta_i)+\lambda_{id}b_d(X_{id};\eta)+\varepsilon_{id},\qquad \lambda_{id}=g(Q_{id},X_{id};\gamma).\]
Here \(b_d\) is a cognitive default, \(g\) maps measurements to shrinkage intensity, and \(\eta,\gamma\) are parameters shared across specified domains. Define prediction loss in a held-out domain \(d\) as \(R_d=\mathbb E[\ell(Y_{id},\widehat Y_{id})]\), for a preregistered loss \(\ell\); estimate the prediction using parameters trained without that domain. Can measured uncertainty and a low-dimensional account of defaults and task difficulty improve \(R_d\) over independently fitted domain models, while preserving calibration across populations and description-versus-experience treatments? Also estimate whether \(Q_{id}\) predicts independently measured response noise. This formulation asks for testable portability and calibration rather than asserting that a universal model exists. Utility benchmarks and scaling of responses must be specified separately before the cross-domain comparison has empirical meaning.

\par\textbf{Formulation and scope.} The source explicitly discusses missing theories, calibration, and additional domains. The unified prediction equation, held-out evaluation, and universal-model wording are our adaptation.

\par\textbf{Open-question provenance.} An explicit question or research agenda was verified at the source locator. This does not by itself certify that every later version remains unresolved \sref{30007583}{1}.

\par\textbf{Historical priority.} The 2023 discussion is the earliest full-text statement verified here. A 2019 working paper exists, but its priority for these particular open extensions was not verified.

\subsection{Short English statement}
Identify whether the specified cognitive-uncertainty model predicts choice biases with portable parameters across economic decision tasks.

\subsection{Sources}
\begin{itemize}\raggedright
\item[\textnormal{[S1]}]\hypertarget{source:30007583:1}{} Benjamin Enke, Thomas Graeber. 2023. Cognitive Uncertainty. Author manuscript dated March 14, 2023, §8 Discussion, printed pp. 33–35 (PDF pp. 34–36)..
\url{https://thomasgraeber.com/assets/papers/EnkeGraeber_CognitiveUncertainty.pdf}.
DOI: 10.1093/qje/qjad025.
{\small\textit{Earliest verified question/agenda reference; historical priority qualified below. The discussion identifies missing general theories of cognitive defaults and task complexity, asks when cognitive uncertainty reflects actual noise, and suggests applications in additional decision domains.}}
\end{itemize}
\end{document}
